\documentclass[main]{subfiles}

\begin{document}
\ifSubfilesClassLoaded{\setcounter{chapter}{4}}{}
\chapter{Gap Analysis and Future Experimental Strategies}

\section{Overview}

Chapter~4 showed that only 19 of the 273 compiled datasets involve an
antimatter sector, and that every matter--antimatter pair currently
matchable from the database is consistent with, but cannot distinguish
from, an ordinary sensitivity gap between the two experiments being
compared. This chapter quantifies exactly where that gap is largest,
by fermion sector (Section~\ref{sec:sector-gaps}) and by DM potential
(Section~\ref{sec:potential-gaps}), identifies data-quality limitations of
the current database that compound the physics gap
(Section~\ref{sec:data-quality}), and proposes where future experimental and
data-compilation effort would do the most to close it
(Section~\ref{sec:strategy}). The guiding principle throughout is the one
established in Chapter~3: because $A_\alpha$ computed from one-sided upper
bounds cannot separate a real CPT signal from a sensitivity gap, the most
valuable next experiment is not necessarily the one that pushes an existing
matter-sector bound tighter, but the one that fills a currently empty
antimatter-sector cell with a bound of \emph{comparable} precision to its
matter-sector counterpart.

\section{Coverage by Fermion Sector}
\label{sec:sector-gaps}

Table~\ref{tab:sector-gaps} compares, for each of the nine canonical
matter--antimatter sector pairs used by SPINDEP, the number of compiled
datasets on each side.

\begin{table}[h]
\centering
\begin{tabular}{@{}l l r r >{\raggedright\arraybackslash}p{3.2cm}@{}}
\toprule
Matter sector & Antimatter sector & Matter & Antimatter & Status \\
\midrule
$e$-$e$     & $e$-$e^+$      & 50 & 5 & Populated (used in Ch.~4) \\
$e$-$p$     & $e$-$\bar p$   & 12 & 9 & Populated (used in Ch.~4) \\
$e$-$\mu$   & $e$-$\bar\mu$  & 3  & 3 & Populated, not yet matched \\
$e$-$n$     & $e$-$\bar n$   & 16 & 0 & \textbf{Gap} \\
$\mu$-$\mu$ & $\mu$-$\bar\mu$& 1  & 0 & \textbf{Gap} (thin on both sides) \\
$n$-$p$     & $n$-$\bar p$   & 9  & 0 & \textbf{Gap} \\
$n$-$n$     & $n$-$\bar n$   & 8  & 0 & \textbf{Gap} \\
$p$-$p$     & $p$-$\bar p$   & 0  & 0 & No data either side \\
$e$-$N$     & $e$-$\bar N$   & 58 & 0 & \textbf{Largest gap} \\
\bottomrule
\end{tabular}
\caption{Matter versus antimatter dataset counts per canonical fermion sector pair.}
\label{tab:sector-gaps}
\end{table}

Table~\ref{tab:sector-gaps}'s nine rows account for 17 of the chapter's 19
antimatter datasets; the remaining two, $\bar p$-$N$ (from antiprotonic
helium) and $\mu^+$-$N$ (from the $dd\mu^+$ molecular ion), have no
matter-side counterpart yet compiled and so fall outside SPINDEP's nine
canonical pairs entirely -- a gap of a different kind from the ones below,
since it is not that the antimatter side is thin, but that there is
currently nothing on the matter side to pair it against at all.

The $e$-$N$ sector is the single largest opportunity in the database: 58
matter-sector datasets exist, more than any other sector, against zero in
$e$-$\bar N$. The $e$-$n$, $n$-$p$, and $n$-$n$ sectors each have a
substantial matter-sector base (8--16 datasets) and no antimatter-sector
counterpart at all. Even the $e$-$\mu$/$e$-$\bar\mu$ pair, nominally
``populated'' on both sides, has not yet been matched into a pair by
SPINDEP's pairing logic (Chapter~4, Section~4.2.3), which is itself worth
investigating: either the compiled $e$-$\bar\mu$ datasets do not share a
common coupling/potential/interaction-class label with the $e$-$\mu$ ones,
or their $\lambda$ ranges do not overlap, and either explanation would be
useful to confirm directly rather than assume.

\section{Coverage by Dobrescu--Mocioiu Potential}
\label{sec:potential-gaps}

Table~\ref{tab:potential-gaps} breaks the same comparison down by DM
potential rather than by sector, restricted to the potentials with at least
one dataset in the database (excluding the 39 \texttt{UNKNOWN}-potential
entries discussed in Section~\ref{sec:data-quality}).

\begin{table}[h]
\centering
\begin{tabular}{@{}l r r >{\raggedright\arraybackslash}p{4.5cm}@{}}
\toprule
Potential & Matter datasets & Antimatter datasets & Status \\
\midrule
$V_{4+5}$  & 58 & 3 & Thin antimatter coverage \\
$V_{9+10}$ & 31 & 0 & \textbf{Gap} \\
$V_{2+3}$  & 21 & 9 & Best-covered potential (used throughout Ch.~4) \\
$V_{1a}$   & 29 & 0 & \textbf{Gap} \\
$V_2$      & 18 & 2 & Thin antimatter coverage \\
$V_1$      & 27 & 3 & Thin antimatter coverage (13 datasets classified by directory or source publication, Ch.~4) \\
$V_{12+13}$& 16 & 0 & \textbf{Gap} \\
$V_{11}$   & 7  & 0 & \textbf{Gap} \\
$V_8$      & 5  & 1 & Thin antimatter coverage \\
$V_{15}$   & 2  & 0 & \textbf{Gap} \\
$V_{16}$   & 2  & 0 & \textbf{Gap} \\
\bottomrule
\end{tabular}
\caption{Matter versus antimatter dataset counts per Dobrescu--Mocioiu potential.}
\label{tab:potential-gaps}
\end{table}

Six of the eleven classified potentials ($V_{11}$, $V_{15}$, $V_{16}$,
$V_{1a}$, $V_{9+10}$, and $V_{12+13}$) have \emph{zero} antimatter-sector
datasets despite substantial matter-sector coverage, $V_{1a}$ and $V_{9+10}$
most strikingly, with 29 and 31 matter-sector datasets respectively and
nothing to compare them against. $V_1$ is the only potential with thin,
rather than zero, antimatter coverage at a comparable matter-sector scale.
Thirteen of its datasets carry no potential token in their filename:
\texttt{Delaunay\_2017} and \texttt{Adkins\_2022\_eeplus} are confirmed
against their source publications directly (Chapter~4,
Section~\ref{sec:database}), and the remaining eleven, including
\texttt{Salumbides\_antipHe\_2014} and \texttt{Salumbides\_ddmu\_2014},
which supply $V_1$'s only antimatter-sector coverage beyond
\texttt{Adkins\_2022\_eeplus}, are resolved from the directory containing
them. Only $V_{2+3}$ has
antimatter coverage exceeding a handful of datasets, which is why nine of
the ten matched pairs in Chapter~4 involve $V_2$ or $V_{2+3}$ (the tenth,
$g_sg_s\cdot V_1\cdot ee$, exists only because of the same
\texttt{Delaunay\_2017}/\texttt{Adkins\_2022\_eeplus} reclassification):
the pairing methodology can only produce a comparison where both sides of
a channel actually have data, and at present that is a narrow slice of the
full sixteen-potential catalogue established in Chapter~2.

\section{Interaction-Range Coverage}

Figures~\ref{fig:lambda-coverage-v1}--\ref{fig:lambda-coverage-v16} show the
dataset inventory for each potential individually (one figure per
potential, rather than a single stitched image, since the full 273-dataset
listing is too tall to render legibly on one page). Coverage is broadly
continuous from sub-picometre to macroscopic scales for the best-populated
potentials ($V_{4+5}$, $V_{9+10}$, $V_{2+3}$), but a small number of
individual datasets report $\lambda$ values many orders of magnitude beyond
the observable universe ($\gtrsim10^{30}$~m for a few entries under
$V_{2+3}$ and $V_{9+10}$), which is physically implausible and almost
certainly reflects a unit-parsing error in the source CSV rather than a
genuine reported bound. Section~\ref{sec:data-quality} returns to this as a
concrete cleanup item rather than a real coverage gap, since including it
uncritically would visibly distort the true $\lambda$-coverage picture.

\begin{figure}[p]
\centering
\includegraphics[width=0.85\textwidth,height=0.9\textheight,keepaspectratio]{lambda_coverage_V1.png}
\caption{Dataset inventory for $V_1$.}
\label{fig:lambda-coverage-v1}
\end{figure}

\begin{figure}[p]
\centering
\includegraphics[width=0.85\textwidth,height=0.9\textheight,keepaspectratio]{lambda_coverage_V1a.png}
\caption{Dataset inventory for $V_{1a}$.}
\label{fig:lambda-coverage-v1a}
\end{figure}

\begin{figure}[p]
\centering
\includegraphics[width=0.85\textwidth,height=0.9\textheight,keepaspectratio]{lambda_coverage_V2.png}
\caption{Dataset inventory for $V_2$.}
\label{fig:lambda-coverage-v2}
\end{figure}

\begin{figure}[p]
\centering
\includegraphics[width=0.85\textwidth,height=0.9\textheight,keepaspectratio]{lambda_coverage_V2p3.png}
\caption{Dataset inventory for $V_{2+3}$.}
\label{fig:lambda-coverage-v23}
\end{figure}

\begin{figure}[p]
\centering
\includegraphics[width=0.85\textwidth,height=0.9\textheight,keepaspectratio]{lambda_coverage_V4p5.png}
\caption{Dataset inventory for $V_{4+5}$, the most populated potential in the database (61 datasets).}
\label{fig:lambda-coverage-v45}
\end{figure}

\begin{figure}[p]
\centering
\includegraphics[width=0.85\textwidth,height=0.9\textheight,keepaspectratio]{lambda_coverage_V8.png}
\caption{Dataset inventory for $V_8$.}
\label{fig:lambda-coverage-v8}
\end{figure}

\begin{figure}[p]
\centering
\includegraphics[width=0.85\textwidth,height=0.9\textheight,keepaspectratio]{lambda_coverage_V9p10.png}
\caption{Dataset inventory for $V_{9+10}$.}
\label{fig:lambda-coverage-v910}
\end{figure}

\begin{figure}[p]
\centering
\includegraphics[width=0.85\textwidth,height=0.9\textheight,keepaspectratio]{lambda_coverage_V11.png}
\caption{Dataset inventory for $V_{11}$.}
\label{fig:lambda-coverage-v11}
\end{figure}

\begin{figure}[p]
\centering
\includegraphics[width=0.85\textwidth,height=0.9\textheight,keepaspectratio]{lambda_coverage_V12p13.png}
\caption{Dataset inventory for $V_{12+13}$.}
\label{fig:lambda-coverage-v1213}
\end{figure}

\begin{figure}[p]
\centering
\includegraphics[width=0.85\textwidth,height=0.9\textheight,keepaspectratio]{lambda_coverage_V15.png}
\caption{Dataset inventory for $V_{15}$.}
\label{fig:lambda-coverage-v15}
\end{figure}

\begin{figure}[p]
\centering
\includegraphics[width=0.85\textwidth,height=0.9\textheight,keepaspectratio]{lambda_coverage_V16.png}
\caption{Dataset inventory for $V_{16}$.}
\label{fig:lambda-coverage-v16}
\end{figure}

\begin{figure}[h]
\centering
\includegraphics[width=0.75\textwidth]{pair_coverage_matrix.png}
\caption{Dataset count by potential and sector; antimatter-sector rows are highlighted, visually reproducing the gaps quantified in Tables~\ref{tab:sector-gaps} and~\ref{tab:potential-gaps}.}
\label{fig:pair-coverage}
\end{figure}

\section{Data-Quality and Framework Limitations}
\label{sec:data-quality}

Three limitations of the current database and framework, distinct from the
physics coverage gaps above, are worth stating plainly rather than working
around silently.

\textbf{Unresolved potential classification.} 39 datasets, roughly a
seventh of the database, carry a filename that cannot be resolved to a
specific $V_n$ token and sit in a directory that does not name the
potential either; these are recorded as \texttt{UNKNOWN} rather than
guessed at. Potential classification otherwise follows three sources of
evidence, in order: an explicit token in the filename; failing that, the
containing directory name, for files whose directory names the potential
unambiguously even though the filename does not (eleven files, including
\texttt{Hoskins\_1985}, \texttt{Kapner\_2007}, \texttt{Chen\_2016},
\texttt{Tan\_2020}, \texttt{Lee\_2020}, \texttt{Bordag\_2001},
\texttt{Alighanbari\_2020}, and \texttt{Colliders}, all resolved to $V_1$
this way); and, where neither is available, direct verification against
the source publication, as for \texttt{Delaunay\_2017} and
\texttt{Adkins\_2022\_eeplus} (also $V_1$; Chapter~4,
Section~\ref{sec:database}), confirmed rather than inferred from the
$g_sg_s$ coupling label alone. Two of the directory-resolved files,
\texttt{Salumbides\_antipHe\_2014} and \texttt{Salumbides\_ddmu\_2014},
carry sector tokens (\texttt{antipHe}, \texttt{ddmu}) registered as
antimatter-sector on the basis of Cong et al.\ (2025)
\cite{Cong2025} (p.~82), which identifies both as $\bar p$-$N$ and
$\mu^+$-$N$ bounds respectively, antimatter sectors despite carrying no
\texttt{bar} suffix in their sector code. The remaining 39 \texttt{UNKNOWN}
datasets have neither a resolvable filename token nor a directory-level
signal; resolving them would require checking each source publication
individually rather than trusting existing metadata, which is left as
future work.

\textbf{Two coupling-label parser artifacts.} The \texttt{Code-plot-matlab}
coupling label affecting 2 datasets (Chapter~4, Section~4.3.1) reflects a
source directory named for a file format rather than a physics category,
and should be reclassified or excluded before any coupling-level summary
statistic is quoted from the database in a publication.

\textbf{Implausible $\lambda$ values in a small number of datasets.} As
noted in Section 5.3, a handful of entries under $V_{2+3}$ and $V_{9+10}$
report interaction ranges that are physically nonsensical
($\lambda\gg10^{26}$~m, larger than the observable universe), almost
certainly the result of a unit token being mis-detected during compilation
(Chapter~4, Section~4.2.2). These entries should be individually
re-verified against their source publications before the database is
considered final; none of the ten matched pairs in Chapter~4, Table~4.2
are affected, since their reported $\lambda_{\max}$ values are all
physically reasonable, but a broader coverage claim drawn from the full
database without this cleanup would not be reliable.

Separately, Chapter~3's $d_{\mu\nu}$ derivation leaves the overall sign of
the momentum-dependent $d_{ij}$ and $d_{00}$ terms open (the mass-enhanced
$d_{i0}\to V_2$ result itself matches the literature exactly); that
limitation is theoretical rather than a database issue, but is relisted
here because it limits how confidently any future $d_{ij}$- or
$d_{00}$-sector experimental target (Section~\ref{sec:strategy}) can be
motivated until it is resolved.

\section{Proposed Experimental Strategies}
\label{sec:strategy}

Three concrete directions follow from the coverage gaps quantified above.

\textbf{Prioritise the $e$-$N$, $n$-$p$, $n$-$n$, and $e$-$n$ sectors for
new antimatter-sector bounds.} These four sectors together account for 91
of the 254 matter-sector datasets in the database and zero antimatter-sector
ones. Precision antiproton experiments such as BASE, and antihydrogen
spectroscopy such as ALPHA, are the most direct existing platforms capable
of producing $\bar p$- or $\bar H$-based bounds in these channels; the
scientific case here is stronger than for further improving already
well-covered channels like $e$-$e^+$, precisely because no comparison is
currently possible at all in these four sectors, regardless of precision.

\textbf{Prioritise $V_{1a}$, $V_{9+10}$, and $V_{12+13}$ for new
antimatter-sector bounds}, for the same reason at the potential level:
these three potentials have 16--31 matter-sector datasets each and none on
the antimatter side, so even a modest first antimatter-sector measurement
would open a channel that currently cannot be compared at all. $V_1$ has
thin but nonzero antimatter coverage (Table~\ref{tab:potential-gaps}) and
so is better placed alongside $V_2$ and $V_8$ as a candidate for
\emph{improving} existing antimatter-side precision rather than opening a
channel from zero.

\textbf{Reinterpret existing antimatter- and Lorentz/CPT-sensitive datasets
before commissioning new experiments.} Cong et al.\ (2025) document a
precedent for this: existing comagnetometer data originally collected to
search for sidereal Lorentz-violation signals has been reinterpreted, after
the fact, as bounds on several $V_n$ terms with the Sun, Moon, or Earth as
an unpolarised source mass. Given that the $g_sg_s\cdot V_1$ pairing in
Chapter~4 already draws on positronium data (Adkins et al., 2022) not
originally collected for this framework, a systematic search of existing BASE, ALPHA, and
comagnetometer publications for data that could be reinterpreted as bounds
in the gap sectors identified in Section~\ref{sec:sector-gaps} is likely to
close some of this gap faster, and at lower cost, than commissioning new
measurements from scratch.

Across all three directions, the standard this thesis applies (Chapter~3)
is that a new antimatter-sector bound is scientifically useful for a CPT
test specifically when it approaches the precision of its matter-sector
counterpart, not merely when it exists: a first antimatter measurement many
orders of magnitude looser than the matter-sector bound in the same channel
would simply reproduce the same sensitivity-gap-dominated
$|A_\alpha|\approx1$ pattern documented throughout Chapter~4, without
resolving the CPT question it is meant to address.

\section{Chapter Summary}

Of the nine canonical fermion-sector pairs and the eleven classified DM
potentials in the compiled database, most have no antimatter-sector
coverage at all: four sectors ($e$-$N$, $e$-$n$, $n$-$p$, $n$-$n$) account
for 91 matter-sector datasets and zero antimatter-sector ones, and six of
eleven potentials are in the same position. This, not a lack of
matter-sector precision, is the binding constraint on this thesis's
ability to distinguish genuine CPT violation from an ordinary sensitivity
gap, and is where future experimental and data-compilation effort would be
most valuable. Chapter~6 concludes the thesis and summarises these findings
alongside those of Chapters~2--4.

\ifSubfilesClassLoaded{
\begin{thebibliography}{99}

\bibitem{Cong2025}
L.~Cong \textit{et al.},
\textit{Spin-dependent exotic interactions},
Rev.\ Mod.\ Phys.\ \textbf{97}, 025005 (2025).

\end{thebibliography}
}{}

\end{document}
